\chapter{极限、连续、完备性与度量空间}
\label{ch:analysis-foundations}

\begin{introduction}
  \item 先修：第 1 章的量词、必要与充分条件、反例和独立 oracle。
  \item 学习产出：能判断一个迭代收敛论证究竟依赖度量、Cauchy 性、完备性还是连续性。
  \item 研究连接：为优化器停止条件、数值不动点、估计量极限和函数逼近建立可审计边界。
\end{introduction}

Python 中“误差越来越小”常被当成收敛证据；数学上还必须回答：误差按什么距离计算，极限是否属于当前空间，结论针对每个点还是同时针对所有点，以及有限精度实验能排除哪些失败模式。本章用压缩迭代作正例，用 $f_n(x)=x^n$ 作反例。

\section{度量把“接近”变成可检验关系}

集合 $X$ 上的函数 $d:X\times X\to[0,\infty)$ 若满足非负与分离、对称及三角不等式，就称为度量，$(X,d)$ 称为\textbf{度量空间}：
\[
d(x,y)=0\Longleftrightarrow x=y,\qquad
d(x,y)=d(y,x),\qquad
d(x,z)\le d(x,y)+d(y,z).
\]
同一对象换一种度量，收敛问题可能改变。有限维向量上的常用范数诱导相同的收敛序列，但停止阈值的数值含义仍不同；研究报告必须同时写出度量和阈值。

\begin{example}
收益向量 $x,y\in\mathbb{R}^T$ 可用 $\lVert x-y\rVert_2$ 衡量整体误差，也可用 $\lVert x-y\rVert_\infty$ 控制最坏时点误差。有限维中恒有 $\lVert x-y\rVert_\infty\le\lVert x-y\rVert_2$，所以二范数不超过 $\varepsilon$ 足以保证每项误差不超过 $\varepsilon$；但“归一化 RMSE 很小”不能推出每项满足同一业务容差，因为其阈值随样本长度缩放。
\end{example}

\section{序列极限、Cauchy 性与完备性}

序列 $(x_n)$ 收敛到 $x$，记作 $x_n\to x$，是指
\[
\forall\varepsilon>0\;\exists N\;\forall n\ge N:\quad d(x_n,x)<\varepsilon.
\]
三角不等式保证极限唯一。若只比较序列内部，则得到 Cauchy 条件：
\[
\forall\varepsilon>0\;\exists N\;\forall m,n\ge N:\quad d(x_m,x_n)<\varepsilon.
\]
每个收敛序列都是 Cauchy 序列；反向需要空间\textbf{完备}。完备性是空间的性质，不是“某个序列看起来稳定”的性质。

\begin{example}
有理数截断序列 $1,1.4,1.41,1.414,\ldots$ 在通常距离下是 Cauchy 序列，却不在 $\mathbb{Q}$ 中收敛，因为其唯一实数极限是 $\sqrt2\notin\mathbb{Q}$。因此 $\mathbb{Q}$ 不完备，而 $\mathbb{R}$ 完备。
\end{example}

\section{连续性连接输入极限与输出极限}

对度量空间 $(X,d_X)$、$(Y,d_Y)$，设 $a$ 是 $X$ 的聚点。当 $x\to a$ 时映射 $T:X\to Y$ 的极限为 $L$，记作 $\lim_{x\to a}T(x)=L$，是指
\[
\forall\varepsilon>0\;\exists\delta>0\;\forall x:\quad
0<d_X(x,a)<\delta\Longrightarrow d_Y(T(x),L)<\varepsilon.
\]
条件 $0<d_X(x,a)$ 表明函数极限只看 $a$ 附近，不依赖 $T(a)$ 的取值，也不要求 $T(a)=L$。若 $a$ 是孤立点，该条件会对任意 $L$ 空真，因而不能定义唯一极限。在度量空间中，聚点处的定义等价于：每个满足 $x_n\to a$ 且 $x_n\ne a$ 的序列都有 $T(x_n)\to L$。

映射 $T$ 在任意 $a\in X$ 连续，是指
\[
\forall\varepsilon>0\;\exists\delta>0:\quad
d_X(x,a)<\delta\Longrightarrow d_Y(T(x),T(a))<\varepsilon.
\]
它等价于下面不排除 $x_n=a$ 的序列判据：
\[
x_n\to a\quad\Longrightarrow\quad T(x_n)\to T(a).
\]
若 $a$ 是聚点，连续性还等价于 $\lim_{x\to a}T(x)=T(a)$；若 $a$ 是孤立点，取足够小的 $\delta$ 即知任意映射都在 $a$ 连续。
连续性允许把极限穿过映射，但不自动给出收敛速度，也不保证算法产生的序列收敛。若实现还包含排序、阈值或裁剪，边界点可能不连续，必须单独分析。

\section{压缩映射为何给出可执行的收敛证明}

设 $(X,d)$ 完备，映射 $T:X\to X$ 满足某个 $0\le q<1$ 及
\[
d(Tx,Ty)\le q\,d(x,y)\qquad(x,y\in X).
\]
从任意 $x_0$ 定义 $x_{n+1}=T(x_n)$。首先逐步使用压缩条件：
\[
d(x_{n+1},x_n)\le q^n d(x_1,x_0).
\]
对 $m>n$，三角不等式与几何级数给出
\[
d(x_m,x_n)
\le\sum_{k=n}^{m-1}d(x_{k+1},x_k)
\le\frac{q^n}{1-q}d(x_1,x_0).
\]
右侧随 $n\to\infty$ 趋于零，所以 $(x_n)$ 是 Cauchy 序列。完备性保证存在 $x^*\in X$ 使 $x_n\to x^*$；压缩条件蕴含连续性，故 $T(x^*)=x^*$。若还有另一个不动点 $y^*$，则
\[
d(x^*,y^*)=d(Tx^*,Ty^*)\le qd(x^*,y^*),
\]
而 $q<1$ 迫使两点距离为零。由此得到存在性、唯一性和误差上界，而不只是几次迭代的截图。

\subsection{固定账本：仿射迭代}

令 $T(x)=0.02+0.8x$、$x_0=0$。不动点由 $x^*=0.02+0.8x^*$ 手算为 $0.1$。五步迭代得到
\[
x_5=0.067232,\qquad |x^*-x_5|=0.032768,
\]
且上面的先验界恰为
\[
\frac{0.8^5}{1-0.8}|x_1-x_0|=0.032768.
\]
这组数值是 notebook 的独立 oracle；程序不得调用被测迭代函数生成 expected。

\section{函数序列：逐点收敛不等于一致收敛}

对函数 $f_n:E\to\mathbb{R}$，逐点收敛允许 $N$ 依赖点 $x$；一致收敛要求同一个 $N$ 同时控制所有 $x\in E$：
\[
\forall\varepsilon>0\;\exists N\;\forall n\ge N\;\forall x\in E:
|f_n(x)-f(x)|<\varepsilon.
\]
在 $E=[0,1]$ 上取 $f_n(x)=x^n$。它逐点收敛到
\[
f(x)=\begin{cases}0,&0\le x<1,\\1,&x=1.\end{cases}
\]
但不一致收敛。对每个 $n$ 选择合法见证点 $x_n=2^{-1/n}<1$，则
\[
|f_n(x_n)-f(x_n)|=x_n^n=\frac12.
\]
因此取 $\varepsilon=0.4$ 时，不存在能同时控制所有 $x$ 的 $N$。见证点依赖 $n$ 正是对一致量词的否定，不是作弊。连续函数的一致极限仍连续，而这里的极限在 1 处不连续，也给出第二条诊断。

\section{独立计算与数值边界}

Notebook 复现 $x_5$、误差界，以及 $n=10,100$ 时的见证误差 $1/2$：
\begin{lstlisting}[language=bash]
uv run python tools/check_learning_unit.py \
  --manifest curriculum/manifest.json --unit upper.ch02
\end{lstlisting}
有限网格很可能漏掉靠近 1 的最坏点；浮点输出也不能证明“对所有 $n$、所有 $x$”。代码在这里是反例见证与算术审计，不是一般证明。

\section*{章末练习}
\begingroup\small
\subsection*{口述概念}
\begin{enumerate}[nosep,leftmargin=*]
  \item 区分序列收敛、Cauchy 性与空间完备性。为什么“迭代内部越来越接近”还不能保证存在允许状态空间内的极限？
  \item 区分函数极限、连续性与一致连续性，并解释函数极限为什么要求目标点是聚点。
  \item 区分函数序列的逐点收敛与一致收敛。为什么反例中的见证点可以随 $n$ 改变？
\end{enumerate}

\subsection*{笔试推导}
\begin{enumerate}[nosep,leftmargin=*]
  \item 使用三角不等式证明度量空间中的序列极限唯一，并标明分离性在哪一步使用。
  \item 从压缩条件推导相邻差、Cauchy 几何级数界、先验误差界和后验误差界；指出完备性与连续性分别在哪里进入。
  \item 证明连续函数的一致极限仍连续；再用该结论或移动见证证明 $f_n(x)=x^n$ 在 $[0,1]$ 上不一致收敛。
\end{enumerate}

\subsection*{数值编程}
\begin{enumerate}[leftmargin=*]
  \item 迭代 $x_{n+1}=0.02+0.8x_n,x_0=0$，显示数值、步长与误差。
  \item 保持固定点为 0.1，把压缩率改为 0.99，比较所需步数。
  \item 对 $x^n$ 比较固定点 $x=1/2$ 与移动点 $x_n=2^{-1/n}$。
\end{enumerate}

\subsection*{研究判断}
\begin{enumerate}[nosep,leftmargin=*]
  \item 优化器连续 20 步满足 $\lVert x_{n+1}-x_n\rVert_2<10^{-6}$。这能否证明收敛到最优解？列出仍缺失的数学与实现证据。
  \item 一个算法的状态被限制为有理数，却逼近 $\sqrt2$。继续增加迭代次数能否解决“不返回合法极限”的问题？给出模型层修复。
  \item 神经网络逼近在一万个验证网格点上误差都小于 $10^{-3}$。研究员据此声称连续输入域上一致误差小于 $10^{-3}$；评价该结论并设计更强证据。
\end{enumerate}
\endgroup

\subsection*{基础衔接：独立计算}
在 $[0,1/2]$ 上证明 $x^n$ 一致趋于零，并给出误差不超过 $10^{-3}$ 所需的最小整数 $n$。

\subsection*{分级提示}
\begingroup\small
\begin{enumerate}[nosep,leftmargin=*]
  \item \textbf{口述概念：}Cauchy 只比较尾部项；完备性把内部稳定转成空间内极限。函数极限穿孔，连续性不穿孔；一致控制要求同一个阈值对全域有效。
  \item \textbf{笔试推导：}唯一性取 $\varepsilon=d(x,y)/3$；压缩证明先把长距离拆成相邻距离的几何级数；一致极限证明插入同一个 $f_N$ 并分成三个 $\varepsilon/3$。
  \item \textbf{数值编程：}先写出本章公式中的目标量，再显示中间计算。改变参数前预测结果方向，运行后说明差异来自模型、抽样还是数值近似。
  \item \textbf{研究判断：}小步长可能是停滞，合法状态空间可能不闭，有限网格可能漏掉移动峰值；分别寻找残差/最优性条件、完备或闭性证据、Lipschitz 或区间界。
\end{enumerate}
\endgroup


\section*{本章小结与 Capstone 连接}

度量定义误差，Cauchy 性描述内部稳定，完备性保证极限仍在空间内，连续性允许传递极限；一致收敛则控制整族输入。上册 Capstone 将用这条链审计迭代算法的停止规则、状态空间和失败见证。
